\documentclass[11pt]{article}
\usepackage[a4paper,margin=1in]{geometry}
\usepackage{amsmath,amsthm,amssymb}
\usepackage{mathtools}

\newtheorem{theorem}{Theorem}[section]
\newtheorem{proposition}[theorem]{Proposition}
\newtheorem{lemma}[theorem]{Lemma}
\newtheorem{corollary}[theorem]{Corollary}
\theoremstyle{definition}
\newtheorem{definition}[theorem]{Definition}
\newtheorem{remark}[theorem]{Remark}
\newtheorem{hypothesis}[theorem]{Hypothesis}

\newcommand{\R}{\mathbb{R}}
\newcommand{\Fr}{\mathcal{A}}
\newcommand{\Om}{\Omega}
\newcommand{\eps}{\varepsilon}
\newcommand{\rstar}{\rho_{*}}
\newcommand{\rhad}{\widehat\rho}
\newcommand{\X}{\mathcal X}

\title{Integrated Action, Kinetic Bound,\\
       and the Weighted Metric Trace Lemma\\
       (Route~$\delta$, Plan~2, Building Block \textsc{WeightedMetricTrace})}
\author{}
\date{}

\begin{document}
\maketitle

\begin{abstract}
This building block chains conditional, explicitly named ingredients into a pointwise metric
distance bound at the deficit endpoint $\rhad:=\rho_\delta=1-\kappa\sqrt{\delta}$,
where throughout this note we write $\delta:=\delta_T(\Om)$.
The repaired version does not use the unresolved centroid
space--time $\Pi$ estimate.  Instead, on the levels where the
isoperimetric defect is small it assumes a same-centre Fusco--Julin
trace package: one bounded Borel centre controls both the Fraenkel
distance and the homothetic trace.  On the same good levels, the low-gradient part of
$\partial^*E_\rho$ is self-truncated and charged by the
Cauchy--Schwarz defect $D_H$ via \cite{TrimmedVelocityRepair}, replacing
the former pointwise lower-gradient hypothesis.  The metric
first-variation and all bad-level kinetic estimates are kept as
conditional inputs, because the current upstream metric-framework
finite-perimeter closure and the old uniform Lipschitz bound (A0) are
not available as proved results.  Under these inputs this yields the integrated action bound
$\int(d_\X(F_\rho,B_1)^2+|\dot F_\rho|_\X^2)\,d\mu \le C\,\delta_T$,
which encodes the integrated isoperimetric defect $D_I$ (via
Fusco--Julin) in its distance term and the integrated Cauchy--Schwarz
defect $D_H$ (via the conditional metric first variation of \cite{MetricFramework})
in its kinetic term.  A Talenti bad-set decomposition then gives the
kinetic bound (K) in Lebesgue form.  The abstract weighted metric trace
lemma combines these two inputs --- Markov on the distance term
(encoding $D_I$) and Cauchy--Schwarz on the kinetic term (encoding
$D_H$) --- on the \emph{order-one} window
$J^*:=[\rho_\delta-L^*,\rho_\delta]$ with $L^*:=(\rho_\delta-\rstar)/4$,
to deliver the \emph{uniform} pointwise bound
$\sup_{\rho\in J^*} d_\X(F_\rho,B_1)^2\le C_*\,\delta_T$.
In particular, setting $\rhad:=\rho_\delta$,
$d_\X(F_{\rho_\delta},B_1)^2\le C_*\,\delta_T$ at the upper endpoint of
the deficit window.  Either input alone would force a $\sqrt{\delta}$
exponent loss; together they yield the sharp bound at $\rho_\delta$.
The chain uses no pointwise (R)/(1.2) and no centroid $\Pi$-closure input.
\end{abstract}

\section{Statements}\label{sec:statements}

Fix $n\ge 2$, $R\ge 2$, $\rstar\in(0,1)$, $\kappa>0$. Let $\Om\subset B_R$
open with $|\Om|=\omega_n$, $u=u_\Om$ the variational torsion function,
$\delta:=\delta_T(\Om)\le\delta_0(n,R,\rstar,\kappa)$.  We always take
\begin{equation}\label{eq:delta-smallness}
   \delta_0\le \left(\frac{1-\rstar}{2\kappa}\right)^2,
\end{equation}
so that
\[
   \rho_\delta:=1-\kappa\sqrt\delta\ge\frac{1+\rstar}{2},
   \qquad
   L^*:=\frac{\rho_\delta-\rstar}{4}\ge\frac{1-\rstar}{8}.
\]
Thus the trace window used below is non-empty and has order-one length.
Set $E_\rho:=\{u>t(\rho)\}$
with $|E_\rho|=\omega_n\rho^n$, $\rho_\delta:=1-\kappa\sqrt\delta$,
$w(\rho):=-t'(\rho)\ge0$, and $\mu(d\rho):=w(\rho)d\rho$. Define
$F_\rho:=\rho^{-1}E_\rho\subset B_{R/\rstar}$, $\gamma(\rho):=[F_\rho]\in\X$.
For a.e.\ $\rho$, $E_\rho$ has finite perimeter: this follows from the
coarea formula applied to $u\in W^{1,2}_0(\Om)$.

\begin{theorem}[Conditional integrated action bound]\label{thm:52int}
\begin{equation}\label{eq:52int}
   \int_{\rstar}^{\rho_\delta}\bigl(d_\X(F_\rho,B_1)^2+|\dot F_\rho|_\X^2\bigr)\,d\mu(\rho)
   \le C(n,R,\rstar,\kappa)\,\delta_T(\Om).
\end{equation}
\end{theorem}

\begin{theorem}[Conditional kinetic bound (K), Lebesgue form]\label{thm:K}
\begin{equation}\label{eq:K}
   \int_{\rstar}^{\rho_\delta}|\dot F_\rho|_\X^2\,d\rho
   \le C'(n,R,\rstar,\kappa)\,\delta_T(\Om).
\end{equation}
\end{theorem}

\begin{theorem}[Weighted metric trace]\label{thm:trace}
Set $L^*:=(\rho_\delta-\rstar)/4$ and
$J^*:=[\rho_\delta-L^*,\rho_\delta]$.  There exists $C_*=C_*(n,R,\rstar,\kappa)$
such that
\begin{equation}\label{eq:trace-uniform}
   \sup_{\rho\in J^*}\,d_\X(F_\rho,B_1)^2\;\le\;C_*\,\delta_T(\Om).
\end{equation}
In particular, setting $\rhad:=\rho_\delta$,
\begin{equation}\label{eq:trace-endpoint}
   d_\X(F_{\rho_\delta},B_1)^2\;\le\;C_*\,\delta_T(\Om),
\end{equation}
i.e.\ the metric distance to the ball is sharply controlled at the
upper endpoint of the deficit window.
\end{theorem}

\subsection{Self-truncated velocity input}\label{ssec:trimmed-velocity}

The metric-derivative bound used below does \emph{not} require a
pointwise lower bound on $|\nabla u|$.  Instead we use the following
levelwise consequence of \cite{TrimmedVelocityRepair}.  On every
Fusco--Julin good level, i.e. on levels where
$D_I(t(\rho))\le\eta_I$, the perimeter is uniformly bounded.  Writing
$f=|\nabla u|$ on $\partial^*E_\rho$ and
$\bar f_\rho=|E_\rho|/P(E_\rho)$, split
\[
   \partial^*E_\rho
   =\{f\ge \bar f_\rho/2\}\cup\{f<\bar f_\rho/2\}.
\]
The low-gradient part satisfies
\[
   \int_{\{f<\bar f_\rho/2\}}\frac1f\,d\mathcal H^{n-1}
   \le C(n,\rstar)\,D_H(t(\rho)),
\]
by the weighted variance identity for $D_H$.  On the complementary part,
$f\ge\bar f_\rho/2$ gives the reciprocal-variance comparison needed in
the old Hyp-G proof.  Combining the two pieces yields
\begin{equation}\label{eq:self-truncated-V}
   \biggl(\int_{\partial^*E_\rho}|V_\rho-\bar V_\rho|\,
     d\mathcal H^{n-1}\biggr)^2
   \le C_{\rm tr}(n,\rstar)\,D_H(t(\rho))
   \qquad(\rho\in G_I).
\end{equation}
Here
\[
   \bar V_\rho:=\frac{1}{P(E_\rho)}
       \int_{\partial^*E_\rho}V_\rho\,d\mathcal H^{n-1}
       =\frac{P(B_\rho)}{P(E_\rho)}.
\]
Thus tiny components or other anomalous low-gradient pieces are not
controlled by a uniform pointwise constant; they are discarded at the
level surface and paid for by $D_H$ itself.

\section{Levelwise quantitative input}\label{sec:FJ}

For a.e.\ $\rho$ write
\[
   \mathcal I(\rho):=D_I(t(\rho)),\qquad
   \mathcal H(\rho):=D_H(t(\rho)),\qquad
   \Delta P_\rho:=P(E_\rho)-P(B_\rho).
\]
Since $P(B_\rho)\ge n\omega_n\rstar^{\,n-1}$ and
$D_I(t(\rho))=(P(E_\rho)-P(B_\rho))(P(E_\rho)+P(B_\rho))$,
\begin{equation}\label{eq:DP-DI}
   0\le \Delta P_\rho\le C_{n,\rstar}\mathcal I(\rho).
\end{equation}
The exact level-set deficit identity of
\cite[Eq.~(eq:main), Thm.~5.1]{LevelSetIdentity}, together with the
normalisation $|\Om|=\omega_n$,
$\delta_T(\Om):=E(\Om)-E(B^*)\ge 0$, reads
\[
   \frac{2\delta_T(\Om)}{\omega_n}
   =\frac{1}{\omega_n}\int_{0}^{\|u\|_\infty}
   \frac{D_H(t)+D_I(t)}{n^2\omega_n^{2/n}m(t)^{1-2/n}}dt.
\]
On the deficit window $[\rstar,\rho_\delta]\subset[0,1]$, the weight
$1/(n^2\omega_n^{2/n}m(t)^{1-2/n})$ is bounded below by
$c_{n,\rstar}>0$ since $m(t)\le \omega_n$.  Therefore
\begin{equation}\label{eq:DH-DI-int}
   \int_{\rstar}^{\rho_\delta}\bigl(\mathcal H(\rho)+\mathcal I(\rho)\bigr)\,
   d\mu(\rho)\le C(n,\rstar)\,\delta_T(\Om),
\end{equation}
where $C(n,\rstar):=2/c_{n,\rstar}=2n^2\omega_n^{2/n+1}\rstar^{n-2}$ for
$n\ge 2$.

\begin{remark}[Sign in the $\rho$ change of variables]\label{rem:sign}
Because the volume radius increases as the threshold decreases,
$t'(\rho)\le0$ and
\[
   dt=t'(\rho)\,d\rho,\qquad d\mu(\rho)=(-t'(\rho))\,d\rho=-dt.
\]
Consequently, for a non-negative level integrand $q$,
\[
   \int_{\rstar}^{\rho_\delta}q(t(\rho))\,d\mu(\rho)
   =\int_{t(\rho_\delta)}^{t(\rstar)}q(t)\,dt.
\]
This is the orientation used in \eqref{eq:DH-DI-int}; the older shorthand
$dt=d\mu$ had the wrong sign.
\end{remark}

\begin{hypothesis}[Same-centre FJ trace package]\label{hyp:FJ-trace}
There are constants $\eta_I>0$, $C=C(n,R,\rstar)$ and a bounded Borel centre
field $\rho\mapsto z_\rho$ such that, on the good set
\[
   G_I:=\{\rho\in[\rstar,\rho_\delta]:\mathcal I(\rho)\le\eta_I\},
\]
the same centre $z_\rho$ satisfies
\begin{align}
   |E_\rho\Delta B_\rho(z_\rho)|^2
      &\le C\,\mathcal I(\rho),\label{eq:FJ-asym}\\
   \left(\int_{\partial^*E_\rho}|H_{z_\rho,\rho}-\bar V_\rho|\,
       d\mathcal H^{n-1}\right)^2
      &\le C\,\mathcal I(\rho),\label{eq:htrace-good}
\end{align}
where
\[
   \bar V_\rho:=\frac{1}{P(E_\rho)}
       \int_{\partial^*E_\rho}V_\rho\,d\mathcal H^{n-1}
       =\frac{P(B_\rho)}{P(E_\rho)}.
\]
This is not imported here from the present FJ bridge: it is the exact
same-centre package still needed upstream.  In particular, a centre that
only minimizes Fraenkel asymmetry, or only minimizes normal oscillation,
is not sufficient for this block.
\end{hypothesis}

\begin{hypothesis}[Good-level metric upper gradient]\label{hyp:metric-good}
For the centres in Hypothesis~\ref{hyp:FJ-trace}, the rescaled curve is
absolutely continuous on $[\rstar,\rho_\delta]$ and, for a.e.\
$\rho\in G_I$,
\[
   |\dot F_\rho|_\X\le C_{n,\rstar}\inf_{a\in\R^n}
   \int_{\partial^*E_\rho}|V_\rho-H_{z_\rho,\rho}-a\cdot\nu_\rho|\,
     d\mathcal H^{n-1}.
\]
This is the finite-perimeter upper-gradient input requested from
\cite{MetricFramework}; the unresolved lower-semicontinuity issue is not
hidden in the proof below.
\end{hypothesis}

\begin{hypothesis}[Bad-level kinetic inputs]\label{hyp:bad-kinetic}
With $B_I:=[\rstar,\rho_\delta]\setminus G_I$ and
\[
   B_{\tau_0}:=\{\rho\in[\rstar,\rho_\delta]:w(\rho)<\tau_0\},
   \qquad \tau_0:=\frac{\rstar}{2n},
\]
assume the two bad-level action bounds
\begin{align}
   \int_{B_I}|\dot F_\rho|_\X^2\,d\mu(\rho)
      &\le C(n,R,\rstar)\,\delta,\label{eq:bad-I-action}\\
   \int_{B_{\tau_0}}|\dot F_\rho|_\X^2\,d\rho
      &\le C(n,R,\rstar)\,\delta.\label{eq:bad-tau-action}
\end{align}
These are the only bad-level kinetic estimates used below.  They replace
the old unsupported uniform Lipschitz assertion (A0).
\end{hypothesis}

\begin{hypothesis}[Profile-gap and Talenti density inputs]\label{hyp:profile-gap}
The density $w=-t'$ satisfies the Talenti upper bound
\[
   0\le w(\rho)\le \frac1n\qquad\text{for a.e. }\rho\in[\rstar,\rho_\delta],
\]
and the profile-gap bad-set estimate
\begin{equation}\label{eq:profile-gap}
   \mathcal L^1(B_{\tau_0})\le C(n,\rstar)\,\delta.
\end{equation}
The estimate \eqref{eq:profile-gap} is the missing integrated
Talenti/profile-gap input from \cite{LevelSetIdentity}.
\end{hypothesis}

\iffalse
Legacy derivation disabled: the audit found that this same-centre FJ
selection/radial-trace proof is not currently justified.  The active
argument above uses Hypothesis~\ref{hyp:FJ-trace} instead.

\begin{lemma}[Fusco--Julin centre package]\label{lem:FJ-package}
There are constants $\eta_I>0$, $C=C(n,R,\rstar)$ and a bounded Borel centre
field $\rho\mapsto z_\rho$ such that, on the good set
\[
   G_I:=\{\rho\in[\rstar,\rho_\delta]:\mathcal I(\rho)\le\eta_I\},
\]
one has
\begin{align}
   |E_\rho\Delta B_\rho(z_\rho)|^2
      &\le C\,\mathcal I(\rho),\label{eq:FJ-asym}\\
   \int_{\partial^*E_\rho}
      \left|\nu_\rho(x)-\frac{x-z_\rho}{|x-z_\rho|}\right|^2
      d\mathcal H^{n-1}(x)
      &\le C\,\mathcal I(\rho).\label{eq:FJ-normal}
\end{align}
\end{lemma}

\begin{proof}
Apply the strong quantitative isoperimetric inequality of Fusco--Julin
\cite{FJ2014} to the rescaled set $\widetilde E:=E_\rho/\rho$, which
has $|\widetilde E|=\omega_n=|B_1|$.  In its joint form \cite[Thm.~1.1]{FJ2014},
there exists an FJ-type centre $\tilde z$ (oscillation-optimal, not
necessarily Fraenkel-optimal) such that
\[
   \alpha(\widetilde E,\tilde z)^2+\beta(\widetilde E,\tilde z)^2
   \le C_n\,\delta_{\rm iso}(\widetilde E),
\]
where $\alpha(\widetilde E,\tilde z)=|\widetilde E\Delta B_1(\tilde z)|/|B_1|$ and
$\beta(\widetilde E,\tilde z)^2:=P(\widetilde E)-(n-1)\int_{\widetilde E}|y-\tilde z|^{-1}dy$.
Translating back via the change of variables $y=x/\rho$, $z:=\rho\tilde z$,
and using $|\widetilde E\Delta B_1(\tilde z)|=\rho^{-n}|E_\rho\Delta B_\rho(z)|$,
$\beta(\widetilde E,\tilde z)^2=\rho^{1-n}\beta(E_\rho,z)^2$, and
$\delta_{\rm iso}(\widetilde E)=\Delta P_\rho/(n\omega_n\rho^{n-1})$, gives
\[
   \biggl(\frac{|E_\rho\Delta B_\rho(z)|}{\rho^n\omega_n}\biggr)^{\!2}
   +\rho^{1-n}\beta(E_\rho,z)^2
   \;\le\;C'_n\,\frac{\Delta P_\rho}{\rho^{n-1}}.
\]
For finite-perimeter sets the divergence identity
\[
   \int_{\partial^*E_\rho}
      \left|\nu_\rho-\frac{x-z}{|x-z|}\right|^2d\mathcal H^{n-1}
   =2\beta(E_\rho,z)^2
\]
follows by testing $\operatorname{div}((x-z)/|x-z|)=(n-1)|x-z|^{-1}$ in
the distributional finite-perimeter divergence theorem
\cite[Thm.~17.6 and Cor.~17.9]{Maggi2012}.  The vector field
$Y(x):=(x-z)/|x-z|$ is bounded ($|Y|=1$ away from $z$) and has locally
integrable divergence on $\R^n$ for $n\ge 2$, so smooth approximations
$\eta_\eps(|x-z|)Y(x)$ with $\eta_\eps\to\mathbf 1$ converge in the
appropriate $BV$-against-finite-perimeter sense; cf.\
\cite[Lemma~1.2]{FJ2014} where this is done explicitly for the
oscillation index.
Since $\rho\in[\rstar,\rho_\delta]$ and $\rho\le 1$, the bounds
\eqref{eq:DP-DI} give \eqref{eq:FJ-asym} (with constant
$C(n,\rstar)\rho^{n+1}\le C(n,\rstar)$) and \eqref{eq:FJ-normal} (with
constant $2C'_n\rho^{n-1}\Delta P_\rho^{-1}\cdot$ etc.\ absorbed).

\emph{Centre localisation.}  If $|z|>R+\rho$, then
$B_\rho(z)\cap B_R=\emptyset$, so $|E_\rho\Delta B_\rho(z)|=|E_\rho|+|B_\rho|
=2\omega_n\rho^n\ge 2\omega_n\rstar^n$.  By \eqref{eq:FJ-asym} this
forces $\mathcal I(\rho)\ge 4\omega_n^2\rstar^{2n}/C(n,\rstar)\ge:\eta_I$,
contradicting $\rho\in G_I$.  Hence on $G_I$, the FJ centre satisfies
$|z_\rho|\le R+1$, a ball depending only on $R$.  Outside $G_I$ set
$z_\rho=0$.

\emph{Borel measurability.}  The map
$(\rho,z)\mapsto|E_\rho\Delta B_\rho(z)|$ is continuous in $z$ for each
$\rho$ (continuity of translation in $L^1$) and Borel in $\rho$ for
each $z$ (coarea), hence Carath\'eodory.  Define the multifunction
$\Gamma(\rho):=\{z\in\overline{B_{R+1}}:|E_\rho\Delta B_\rho(z)|\le
\inf_{z'\in\overline{B_{R+1}}}|E_\rho\Delta B_\rho(z')|+\rho^n\eta_I/2\}$
on $G_I$.  For each $\rho\in G_I$, $\Gamma(\rho)$ is non-empty
(infimum attained in the compact set $\overline{B_{R+1}}$ by
continuity), closed (preimage of $(-\infty,\inf+\rho^n\eta_I/2]$ under
continuous $z\mapsto|E_\rho\Delta B_\rho(z)|$), and contained in the
fixed compact set $\overline{B_{R+1}}$ by the centre-localisation
step above.  The Kuratowski--Ryll-Nardzewski / Castaing measurable
selection theorem \cite[Thm.~III.6]{CastaingValadier1977} applies and
yields a Borel selector $\rho\mapsto z_\rho$ on $G_I$.
\end{proof}

\begin{lemma}[Radial trace from oriented divergence]\label{lem:radial-slicing}
For the centres of Lemma~\ref{lem:FJ-package} and every $\rho\in G_I$,
\begin{equation}\label{eq:radial-slicing}
   T_\rho:=\int_{\partial^*E_\rho}
      \left|\frac{|x-z_\rho|}{\rho}-1\right|\,d\mathcal H^{n-1}(x)
   \le C(n,R,\rstar)\,\mathcal I(\rho)^{1/2}.
\end{equation}
\end{lemma}

\begin{proof}
Translate so that $z_\rho=0$ and write $e_r=x/|x|$.  Set
$w(r):=|r/\rho-1|$ and \(X(x):=w(|x|)e_r\).  Since \(E_\rho\subset B_R\)
and \(z_\rho\) is bounded, \(w(|x|)\le C(R,\rstar)\) on
\(\partial^*E_\rho\).  Up to the irrelevant point \(x=0\)
(when $z_\rho\in\overline{E_\rho}\cap\partial^*E_\rho$, the truncation
argument below still applies since the singularity at $x=z_\rho$ is a
single-point set, which is $\mathcal H^{n-1}$-null in dimension $n\ge2$),
\[
   1=e_r\cdot\nu_\rho+(1-e_r\cdot\nu_\rho),
\]
therefore
\[
   T_\rho
   =\int_{\partial^*E_\rho}X\cdot\nu_\rho\,d\mathcal H^{n-1}
    +\int_{\partial^*E_\rho}w(|x|)(1-e_r\cdot\nu_\rho)\,d\mathcal H^{n-1}.
\]
The second term is controlled by the normal oscillation, since
\[
   1-e_r\cdot\nu_\rho=\frac12|e_r-\nu_\rho|^2.
\]
For the first term use the finite-perimeter divergence theorem, first with
the truncated smooth fields
\[
   X_{\eps,M}(x)=\eta_\eps(|x|)\,\chi_M(|x|)\,w(|x|)\,e_r
\]
and then let \(\eps\downarrow0\), \(M\uparrow\infty\).  The truncation is
legitimate because \(X\) is bounded on \(B_R\), its divergence is locally
integrable for \(n\ge2\), and \(E_\rho\) has finite perimeter.  Since
\(w(\rho)=0\), the same boundary integral over \(B_\rho\) is zero, hence
\[
   \int_{\partial^*E_\rho}X\cdot\nu_\rho\,d\mathcal H^{n-1}
   =\int_{E_\rho\setminus B_\rho}\operatorname{div}X\,dx
     -\int_{B_\rho\setminus E_\rho}\operatorname{div}X\,dx .
\]
For \(r>\rho\),
\[
   \operatorname{div}X=\frac n\rho-\frac{n-1}{r}\le \frac n\rho,
\]
while for \(0<r<\rho\),
\[
   \operatorname{div}X=\frac{n-1}{r}-\frac n\rho,
   \qquad
   \bigl(-\operatorname{div}X\bigr)_+\le\frac n\rho .
\]
Thus
\[
   \int_{\partial^*E_\rho}X\cdot\nu_\rho\,d\mathcal H^{n-1}
   \le C(n,\rstar)|E_\rho\Delta B_\rho|.
\]
Combining the two estimates gives
\[
   T_\rho\le C(n,R,\rstar)\left(
      |E_\rho\Delta B_\rho(z_\rho)|
      +\int_{\partial^*E_\rho}
        \left|\nu_\rho-\frac{x-z_\rho}{|x-z_\rho|}\right|^2
        d\mathcal H^{n-1}\right).
\]
Finally use \eqref{eq:FJ-asym} and \eqref{eq:FJ-normal}.  Since
\(\mathcal I(\rho)\le\eta_I\le1\) on \(G_I\), this yields
\eqref{eq:radial-slicing}.
\end{proof}

\begin{lemma}[Good-level homothetic trace]\label{lem:good-htrace}
Let
\[
   \bar V_\rho:=\frac{1}{P(E_\rho)}
       \int_{\partial^*E_\rho}V_\rho\,d\mathcal H^{n-1}
       =\frac{P(B_\rho)}{P(E_\rho)}.
\]
For every $\rho\in G_I$,
\begin{equation}\label{eq:htrace-good}
   \left(\int_{\partial^*E_\rho}|H_{z_\rho,\rho}-\bar V_\rho|\,
       d\mathcal H^{n-1}\right)^2
   \le C(n,R,\rstar)\,\mathcal I(\rho).
\end{equation}
\end{lemma}

\begin{proof}
With $e_\rho(x):=(x-z_\rho)/|x-z_\rho|$,
\[
   |H_{z_\rho,\rho}-1|
   \le \left|\frac{|x-z_\rho|}{\rho}-1\right|
      + |e_\rho\cdot\nu_\rho-1|.
\]
Since $|e_\rho|=|\nu_\rho|=1$,
$|e_\rho\cdot\nu_\rho-1|=\frac12|e_\rho-\nu_\rho|^2$.  Hence
\eqref{eq:FJ-normal} and Lemma~\ref{lem:radial-slicing} imply
$\int|H_{z_\rho,\rho}-1|\le C\mathcal I(\rho)^{1/2}$ on $G_I$.  Moreover
\[
   P(E_\rho)|1-\bar V_\rho|=P(E_\rho)-P(B_\rho)=\Delta P_\rho
   \le C\mathcal I(\rho)\le C\mathcal I(\rho)^{1/2}
\]
since $\mathcal I(\rho)\le\eta_I\le 1$ on $G_I$.  The triangle
inequality gives \eqref{eq:htrace-good}.
\end{proof}

\fi

\section{The integrated action bound (5.2-int)}\label{sec:52int}

\begin{proof}[Proof of Theorem~\ref{thm:52int}]
Let $B_I:=[\rstar,\rho_\delta]\setminus G_I$.  By \eqref{eq:DH-DI-int},
\begin{equation}\label{eq:BI-measure}
   \mu(B_I)\le \eta_I^{-1}\int_{\rstar}^{\rho_\delta}\mathcal I(\rho)\,d\mu(\rho)
   \le C\,\delta_T(\Om).
\end{equation}

\paragraph{Distance term.}
On $G_I$, the quotient metric and the rescaling give
\begin{equation}\label{eq:dist-trivial}
   d_\X(F_\rho,B_1)
   \le \rho^{-n}|E_\rho\Delta B_\rho(z_\rho)|,
\end{equation}
hence \eqref{eq:FJ-asym} implies
$d_\X(F_\rho,B_1)^2\le C\mathcal I(\rho)$.  On $B_I$ we use only the
trivial bound $d_\X(F_\rho,B_1)\le 2\omega_n$, because both sets have
volume $\omega_n$.  Therefore, by \eqref{eq:DH-DI-int} and
\eqref{eq:BI-measure},
\begin{equation}\label{eq:dist-action}
   \int_{\rstar}^{\rho_\delta}d_\X(F_\rho,B_1)^2\,d\mu(\rho)
   \le C\,\delta_T(\Om).
\end{equation}

\paragraph{Metric-derivative term.}
On $G_I$, Hypothesis~\ref{hyp:metric-good} gives the per-$\rho$ velocity bound
\[
   |\dot F_\rho|_\X\;\le\;C_{n,\rstar}\inf_{a\in\R^n}\!
   \int_{\partial^*E_\rho}|V_\rho-H_{z_\rho,\rho}-a\cdot\nu_\rho|\,
     d\mathcal H^{n-1}.
\]
We specialise to $a=0$ (a legitimate upper bound), then apply the
triangle inequality
$|V_\rho-H_{z_\rho,\rho}|\le|V_\rho-\bar V_\rho|+|H_{z_\rho,\rho}-\bar V_\rho|$
under the integral, followed by $(a+b)^2\le 2a^2+2b^2$, to obtain
\begin{equation}\label{eq:Fdot-decomp}
   |\dot F_\rho|_\X^{\,2}\le 2C_{n,\rstar}^{2}
   \left[
      \biggl(\int_{\partial^*E_\rho}|V_\rho-\bar V_\rho|\biggr)^{\!2}
      +
      \biggl(\int_{\partial^*E_\rho}|H_{z_\rho,\rho}-\bar V_\rho|\biggr)^{\!2}
   \right].
\end{equation}
The first square is controlled by the self-truncated velocity estimate
\eqref{eq:self-truncated-V} from \cite{TrimmedVelocityRepair}:
\begin{equation}\label{eq:V-bar-V}
   \biggl(\int_{\partial^*E_\rho}|V_\rho-\bar V_\rho|\,
     d\mathcal H^{n-1}\biggr)^{\!2}
   \;\le\;C_{\rm tr}(n,\rstar)\,\mathcal H(\rho),
\end{equation}
and the second square is \eqref{eq:htrace-good}.  Thus
\[
   |\dot F_\rho|_\X^2\le C(n,R,\rstar)\bigl(\mathcal H(\rho)+\mathcal I(\rho)\bigr)
   \qquad\text{for a.e. }\rho\in G_I,
\]
where $C(n,R,\rstar)$ is polynomial in the constants of the
self-truncated velocity estimate and of the Fusco--Julin package.
On $B_I$ we use only the conditional bad-isoperimetric-level action
input \eqref{eq:bad-I-action}.  Combining this with
\eqref{eq:DH-DI-int} gives
\begin{equation}\label{eq:kinetic-action}
   \int_{\rstar}^{\rho_\delta}|\dot F_\rho|_\X^2\,d\mu(\rho)
   \le C\,\delta_T(\Om).
\end{equation}
Equations \eqref{eq:dist-action} and \eqref{eq:kinetic-action} prove
\eqref{eq:52int}.
\end{proof}

\begin{remark}
The good-level estimate is genuinely quantitative: the only square-root
loss is the Fraenkel term from the same-centre FJ input, and it is squared
before integration.  The large-isoperimetric-defect levels are not
estimated pointwise in this file; their distance contribution is paid for
by $\mu(B_I)\le C\delta$, while their kinetic contribution is exactly the
conditional input \eqref{eq:bad-I-action}.
\end{remark}

\section{The kinetic bound (K)}\label{sec:K}

\begin{remark}[No use of the old A0]\label{rem:no-A0}
The previous draft asserted a perimeter-free uniform metric Lipschitz
bound (A0) from nestedness and bounded support.  The audit found this
unsupported.  The Lebesgue kinetic estimate below therefore uses only
the weighted action already proved and the explicit bad-density kinetic
input \eqref{eq:bad-tau-action}.
\end{remark}

\begin{proof}[Proof of Theorem~\ref{thm:K}]
Split $[\rstar,\rho_\delta]=G_{\tau_0}\sqcup B_{\tau_0}$ with
$G_{\tau_0}:=\{\rho:w(\rho)\ge\tau_0\}$ and
$\tau_0=\rstar/(2n)$.  On $G_{\tau_0}$,
\[
   \int_{G_{\tau_0}}|\dot F_\rho|_\X^2\,d\rho
   \le\tau_0^{-1}\int_{G_{\tau_0}}|\dot F_\rho|_\X^2\,d\mu(\rho)
   \le C(n,\rstar)\delta
\]
by Theorem~\ref{thm:52int}.  On $B_{\tau_0}$ use the conditional
bad-density action input \eqref{eq:bad-tau-action}.
Adding gives \eqref{eq:K}.
\end{proof}

\section{The weighted metric trace lemma}\label{sec:agent4}

\begin{theorem}[Abstract weighted metric trace]\label{thm:abstract-trace}
Let $\mu=w(\rho)d\rho$ on $[\rstar,\rho_\delta]$.  Assume:
\[
\begin{array}{ll}
\mathrm{(H1)}&
   \displaystyle \int_{\rstar}^{\rho_\delta}
      \bigl(d_\X(\gamma(\rho),B_1)^2+|\dot\gamma|^2\bigr)\,d\mu
      \le C\delta,\\[2mm]
\mathrm{(H2)}&
   \displaystyle \mu([a,b])\ge c(b-a)-C_{\rm low}\delta
      \quad\text{for every }[a,b]\subset[\rstar,\rho_\delta],\\[2mm]
\mathrm{(H3)}&
   \displaystyle 0\le w(\rho)\le M
      \quad\text{for a.e. }\rho,\\[2mm]
\mathrm{(H4)}&
   \displaystyle \int_{\rstar}^{\rho_\delta}|\dot\gamma|^2\,d\rho
      \le C'\delta .
\end{array}
\]
Set
$L^*:=(\rho_\delta-\rstar)/4$ and $J^*:=[\rho_\delta-L^*,\rho_\delta]$.
Then
\[
   \sup_{\rho\in J^*}\,d_\X(\gamma(\rho),B_1)^2\;\le\;C_*\,\delta,
\]
provided $\delta\le cL^*/(2C_{\rm low})$.  Here
$C_*:=2K_1+2L^*C'$ with $K_1:=4C/(cL^*)$. In particular at the
endpoint, $d_\X(\gamma(\rho_\delta),B_1)^2\le C_*\delta$.
\end{theorem}

\begin{proof}
\emph{Step 1: macroscopic Markov on the order-one window
$J^*=[\rho_\delta-L^*,\rho_\delta]$.}
By (H2), $\mu(J^*)\ge cL^*-C_{\rm low}\delta\ge cL^*/2$. By
(H1) and Chebyshev (Markov on the square),
\[
   \mu\bigl(\{\rho\in J^*:d_\X(\gamma(\rho),B_1)^2>K_1\delta\}\bigr)
   \;\le\;\frac{C\delta}{K_1\delta}\;<\;\mu(J^*)
\]
for $K_1:=4C/(cL^*)$ (any choice with $K_1>2C/\mu(J^*)$ works). Hence
there exists $\rho_0\in J^*$ with
$d_\X(\gamma(\rho_0),B_1)^2\le K_1\delta$.

\emph{Step 2: kinetic transport across $J^*$.} For any $\rho\in J^*$,
$|\rho-\rho_0|\le L^*$ (the diameter of $J^*$). AGS Thm 1.1.2
\cite[Thm.~1.1.2]{AGS2008} gives
$d_\X(\gamma(\rho_0),\gamma(\rho))\le\int_{\rho_0}^{\rho}|\dot\gamma(s)|\,ds$;
Cauchy--Schwarz and (H4) then yield
\[
   d_\X(\gamma(\rho_0),\gamma(\rho))^2
   \;\le\;|\rho-\rho_0|\!\!\int_{\rho_0}^{\rho}|\dot\gamma|^2\,ds
   \;\le\;L^*\!\cdot C'\delta.
\]

\emph{Step 3: triangle inequality.} Squaring
$d_\X(\gamma(\rho),B_1)\le d_\X(\gamma(\rho),\gamma(\rho_0))+d_\X(\gamma(\rho_0),B_1)$
with $(a+b)^2\le 2a^2+2b^2$ gives, uniformly for $\rho\in J^*$,
\[
   d_\X(\gamma(\rho),B_1)^2\;\le\;2L^*C'\delta+2K_1\delta\;=\;C_*\delta.
\]
The endpoint case $\rho=\rho_\delta$ is included since $\rho_\delta\in J^*$.
\end{proof}

\begin{remark}[Why two inputs are needed: the $D_H/D_I$ pairing]
\label{rem:why-two-inputs}
The proof is the 1D Sobolev embedding
$H^1(J^*)\hookrightarrow L^\infty(J^*)$ applied to
$\rho\mapsto d_\X(\gamma(\rho),B_1)$, decomposed into its two ingredients:

(i)~the \emph{distance term} (H1, $L^2(d\mu)$) encodes the integrated
isoperimetric defect $D_I$ via the Fusco--Julin estimate
\eqref{eq:FJ-asym}: many levels are close to balls;

(ii)~the \emph{kinetic term} (H4, $L^2(d\rho)$) encodes the
integrated Cauchy--Schwarz defect $D_H$ via the metric first variation
of \cite{MetricFramework}: the family $\rho\mapsto F_\rho$ does not
``curve away'' from its trajectory faster than $\sqrt\delta$ in
$L^2$.

Step 1 (Markov) uses only (i) and extracts a single $\rho_0$ where
$d^2\le K_1\delta$; it cannot place $\rho_0$ at $\rho_\delta$.
Step 2 (kinetic Cauchy--Schwarz) uses only (ii) and propagates the
bound from $\rho_0$ to any $\rho\in J^*$ at cost
$|\rho-\rho_0|\cdot C'\delta\le L^*C'\delta$.

\emph{Either input alone fails.} If one tried to extract the endpoint
bound by Markov on a window of size $\epsilon$ near $\rho_\delta$
(to keep the boundary-layer transfer cost in
\cite{BoundaryLayerTransfer} of order $\epsilon$), the Markov constant
becomes $4C/(c\epsilon)$, forcing the bound
$d^2(\gamma(\rho_\delta),B_1)\lesssim\delta/\epsilon+\epsilon^2$.
Optimising $\epsilon\sim\delta^{1/3}$ yields the lossy exponent
$\delta^{2/3}$ instead of $\delta$.  Conversely, kinetic Cauchy--Schwarz
alone gives only $\sqrt\delta$-H\"older regularity of $\rho\mapsto d^2$,
which without the integrated $L^2(d\mu)$ control of the distance term
is vacuous as a pointwise bound.

The combination on the order-one window $J^*=[\rho_\delta-L^*,\rho_\delta]$
yields the sharp $\delta$ exponent at the endpoint $\rho_\delta$,
which is exactly the upper boundary of the deficit interval and the
location where the boundary-layer transfer of
\cite{BoundaryLayerTransfer} has thickness $1-\rho_\delta\asymp\sqrt\delta$.
\end{remark}

\paragraph{Verification of hypotheses for the torsion application:}
\begin{itemize}
\item (H1) is Theorem~\ref{thm:52int}.
\item (H2) follows from Hypothesis~\ref{hyp:profile-gap}: on
$G_{\tau_0}:=\{\rho:w(\rho)\ge\tau_0\}$, with
$\tau_0=\rstar/(2n)$, we have $\mu(d\rho)\ge\tau_0\,d\rho$, while
$\mathcal L^1(B_{\tau_0})\le C\delta$.  Hence
\[
   \mu([a,b])\ge \tau_0(b-a)-\tau_0\,\mathcal L^1(B_{\tau_0})
   \ge \tau_0(b-a)-C_{\rm low}\delta.
\]
\item (H3) is the Talenti upper-density part of
Hypothesis~\ref{hyp:profile-gap}, with $M=1/n$.
\item (H4) is Theorem~\ref{thm:K}.
\end{itemize}

\begin{proof}[Proof of Theorem~\ref{thm:trace}]
After fixing the constants in (H2), shrink $\delta_0$ if necessary so
that the smallness condition in Theorem~\ref{thm:abstract-trace} holds
on the order-one window guaranteed by \eqref{eq:delta-smallness}.  Apply
Theorem~\ref{thm:abstract-trace} to $\gamma=[F_\rho]$.
\end{proof}

\begin{corollary}[Pointwise asymmetry at $\rho_\delta$]\label{cor:pointwise-asymmetry}
Setting $\rhad:=\rho_\delta$,
\[
   |E_{\rho_\delta}\Delta B_{\rho_\delta}|^2 \;\le\; C(n,\rstar,R)\,\delta_T(\Om).
\]
\end{corollary}

\begin{proof}
By the rescaling $F_\rho=\rho^{-1}E_\rho$ recorded at
\eqref{eq:dist-trivial}, the endpoint bound \eqref{eq:trace-endpoint}
yields
\[
   |E_{\rho_\delta}\Delta B_{\rho_\delta}|^2
   \le C(n,\rstar,R)\,\rho_\delta^{\,2n}
      d_\X(F_{\rho_\delta},B_1)^2
   \le C(n,\rstar,R)\,\delta_T(\Om).
\]
The translation realising the infimum in $d_\X$ provides the centring
required by the Fraenkel asymmetry on the right-hand side.
\end{proof}

\section{Conditionality}\label{sec:cond}

\begin{itemize}
\item Theorem~\ref{thm:52int}: consumes the exact level-set deficit
budget \eqref{eq:DH-DI-int} from \cite{LevelSetIdentity}, the
same-centre FJ trace package Hypothesis~\ref{hyp:FJ-trace}, the
good-level metric upper gradient Hypothesis~\ref{hyp:metric-good}, the
self-truncated velocity estimate of \cite{TrimmedVelocityRepair} on
$G_I$ only, and the bad-isoperimetric-level kinetic input
\eqref{eq:bad-I-action}.
\item Theorem~\ref{thm:K}: consumes Theorem~\ref{thm:52int} on
$G_{\tau_0}$ and the bad-density kinetic input
\eqref{eq:bad-tau-action}.  It does not use the old A0 assertion.
\item Theorem~\ref{thm:trace}: consumes the cleaned hypotheses
(H1)--(H4).  In the torsion application, (H2) and (H3) are supplied by
Hypothesis~\ref{hyp:profile-gap}.
\end{itemize}

Thus this block no longer depends on the centroid integrated
\(\Pi\)-control input isolated in \cite{SpaceTimeIdentity}, and it no
longer depends on the pointwise lower-gradient hypothesis Hyp-G.  The
low-gradient part of an isoperimetric-good level surface is handled by
truncation and paid for by \(D_H\).  Pointwise (R)/(1.2) is never
invoked as a proved input; the unresolved same-centre FJ/homothetic trace
bridge is isolated in Hypothesis~\ref{hyp:FJ-trace}.

\begin{thebibliography}{99}

\bibitem{LevelSetIdentity}
Building block \textsc{LevelSetIdentity}.

\bibitem{MetricFramework}
Building block \textsc{MetricFramework}.

\bibitem{SpaceTimeIdentity}
Building block \textsc{SpaceTimeIdentity}.

\bibitem{BoundaryLayerTransfer}
Building block \textsc{BoundaryLayerTransfer}.

\bibitem{TrimmedVelocityRepair}
Building block \textsc{TrimmedVelocityRepair}.
Self-truncated velocity defect replacing Hyp-G(lower).

\bibitem{AGS2008}
L.~Ambrosio, N.~Gigli, G.~Savar\'e, \emph{Gradient Flows in Metric
Spaces}, Birkh\"auser, 2008. Thm 1.1.2.

\bibitem{Talenti1976}
G.~Talenti, Ann. Mat. Pura Appl. (4) \textbf{110} (1976), 353--372.

\bibitem{FJ2014}
N.~Fusco, V.~Julin, Calc. Var. PDE \textbf{50} (2014), 925--937.

\bibitem{Maggi2012}
F.~Maggi, \emph{Sets of Finite Perimeter and Geometric Variational
Problems}, Cambridge University Press, 2012.

\bibitem{FigalliMaggi2011}
A.~Figalli, F.~Maggi, Arch. Ration. Mech. Anal. \textbf{201} (2011),
143--207.

\bibitem{AFP2000}
L.~Ambrosio, N.~Fusco, D.~Pallara, \emph{Functions of Bounded Variation
and Free Discontinuity Problems}, Oxford University Press, 2000.

\bibitem{GlobalAssembly}
Building block \textsc{GlobalAssembly}.

\bibitem{GilbargTrudinger2001}
D.~Gilbarg, N.~S.~Trudinger,
\emph{Elliptic Partial Differential Equations of Second Order},
Classics in Mathematics, Springer-Verlag, Berlin (2001 reprint).

\bibitem{BrothersZiemer1988}
J.~E.~Brothers, W.~P.~Ziemer,
\emph{Minimal rearrangements of Sobolev functions},
J.\ Reine Angew.\ Math.\ \textbf{384} (1988), 153--179.

\bibitem{CastaingValadier1977}
C.~Castaing, M.~Valadier,
\emph{Convex Analysis and Measurable Multifunctions},
Lecture Notes in Math.\ \textbf{580}, Springer, 1977.

\bibitem{CianchiMaz'ya2011}
A.~Cianchi, V.~G.~Maz'ya,
\emph{Global Lipschitz regularity for a class of quasilinear
elliptic equations}, J.~Eur.~Math.~Soc.\ \textbf{13} (2011), 247--292.

\end{thebibliography}

\end{document}
